% Secuencia del pipeline: lectura en la subida, trabajo asíncrono por cola, resolución y optimización. Misma paleta que c4-contenedores.tex.
\definecolor{tkPerF}{HTML}{FFF0B8}\definecolor{tkPerB}{HTML}{9A7B00}
\definecolor{tkAppF}{HTML}{CFE2F7}\definecolor{tkAppB}{HTML}{1F4E86}
\definecolor{tkProcF}{HTML}{D3EDD0}\definecolor{tkProcB}{HTML}{2E6B2B}
\definecolor{tkExtF}{HTML}{FAD7B5}\definecolor{tkExtB}{HTML}{A8500F}
\definecolor{tkInfF}{HTML}{DEDEDE}\definecolor{tkInfB}{HTML}{4D4D4D}
\definecolor{tkBdF}{HTML}{CDEBEB}\definecolor{tkBdB}{HTML}{1C6E6E}
\definecolor{tkMqF}{HTML}{E2D6F4}\definecolor{tkMqB}{HTML}{5B2E99}
\begin{tikzpicture}[x=1cm, y=1cm, font=\small,
  pa/.style={draw=tkAppB, fill=tkAppF, rounded corners=2pt, align=center, minimum width=1.95cm, minimum height=9mm},
  pg/.style={draw=tkInfB, fill=tkInfF, rounded corners=2pt, align=center, minimum width=1.95cm, minimum height=9mm},
  pp/.style={draw=tkProcB, fill=tkProcF, rounded corners=2pt, align=center, minimum width=1.95cm, minimum height=9mm},
  pc/.style={draw=tkBdB, fill=tkBdF, rounded corners=2pt, align=center, minimum width=1.95cm, minimum height=9mm},
  pr/.style={draw=tkPerB, fill=tkPerF, rounded corners=2pt, align=center, minimum width=1.95cm, minimum height=9mm},
  pe/.style={draw=tkExtB, fill=tkExtF, rounded corners=2pt, align=center, minimum width=1.95cm, minimum height=9mm},
  pm/.style={draw=tkMqB, fill=tkMqF, rounded corners=2pt, align=center, minimum width=1.95cm, minimum height=9mm},
  life/.style={dashed, draw=tkInfB!70},
  msg/.style={-{Stealth[length=1.8mm]}, thick, draw=tkAppB},
  ret/.style={-{Stealth[length=1.8mm]}, thick, dashed, draw=tkInfB},
  lab/.style={fill=white, inner sep=1pt, align=center}]
\node[pa] (w) at (0,0) {Acudiente\\(web)};
\node[pg] (g) at (2.2,0) {Puerta de\\enlace};
\node[pp] (i) at (4.4,0) {Inteli-\\gencia};
\node[pe] (m) at (6.6,0) {Gemini};
\node[pm] (q) at (8.8,0) {Cola y\\trabajador};
\node[pc] (c) at (11.0,0) {Catálogo};
\node[pr] (r) at (13.2,0) {Revisión\\humana};
\foreach \n in {w,g,i,m,q,c,r} \draw[life] (\n.south) -- ++(0,-10.7);
\draw[msg] (0,-1.1) -- (2.2,-1.1) node[midway, above=1pt, lab] {1 sube};
\draw[msg] (2.2,-1.8) -- (4.4,-1.8) node[midway, above=1pt, lab] {2 límite 429};
\draw[msg] (4.4,-2.5) -- (6.6,-2.5) node[midway, above=1pt, lab] {3 lee páginas};
\draw[ret] (6.6,-3.1) -- (4.4,-3.1) node[midway, above=1pt, lab] {4 renglones};
\draw[msg] (4.4,-3.8) -- (8.8,-3.8) node[midway, above=1pt, lab] {5 guarda el trabajo y lo encola};
\draw[ret] (4.4,-4.5) -- (0,-4.5) node[pos=0.5, above=1pt, lab] {6 responde 202: trabajo aceptado};
\draw[msg] (8.8,-5.3) -- ++(0.8,0) -- ++(0,-0.55) -- ++(-0.8,0);
\node[anchor=west, lab, align=left] at (9.7,-5.55) {7 toma el trabajo\\(RabbitMQ o tabla)};
\draw[msg] (8.8,-6.5) -- (11.0,-6.5) node[midway, above=1pt, lab] {8 pide candidatos};
\draw[ret] (11.0,-7.1) -- (8.8,-7.1) node[midway, above=1pt, lab] {9 candidatos};
\draw[msg] (8.8,-7.8) -- ++(0.8,0) -- ++(0,-1.0) -- ++(-0.8,0);
\node[anchor=west, lab, align=left] at (9.7,-8.3) {10 clasifica: sin candidatos, hueco;\\parecido $\geq \tau = 0{,}75$, resuelto;\\si no, ambiguo};
\draw[msg] (8.8,-9.4) -- ++(0.8,0) -- ++(0,-0.55) -- ++(-0.8,0);
\node[anchor=west, lab, align=left] at (9.7,-9.65) {11 NSGA-II sobre los renglones\\con candidatos};
\draw[msg] (8.8,-10.6) -- (13.2,-10.6) node[midway, above=1pt, lab] {12 kits por revisar};
\end{tikzpicture}
